\section{Discussion and conclusion}
The count measurement, read only after decoding failure, carries a selection effect. Under sign imbalance this effect persists after standardization.
Its strength is governed by the closed-form constant $\eta_k\rho^2$, and its direction reverses between $k=4$ and $k=5$; improving the unconditional
reference distribution does not remove it. For the actual protocol, when $k\ge5$, the one-sided imbalance places the bias on the side of insufficient capacity;
the balanced imbalance is unaffected by selection. The observable sign totals suffice to construct an upper confidence bound after failure,
and both implementation rules, the envelope rule and the mixture rule, pass the end-to-end checks.

Along the load axis, the failure mechanism passes gradually from the macroscopic core of the critical window to edge repetition at a fixed subcritical load, and the failure-conditioned distribution is an explicit two-layer mixture.
When $k\ge5$, the mechanism that dominates the failure probability separates from the mechanism that dominates the conditional mean: the layer where the mean changes sign lies at
$t^2\approx(2k-3)\log M$, deeper than the equal-probability layer. In finite size, the threshold center is given by
the Dembo--Montanari shift, and the shape of the unconditional measurement is given by the chi-square
$(M-1)$; neither of the two involves a fitted parameter.

This paper's theory covers fixed $k\ge3$, the ideal Plain mapping, and deterministic signs independent of the mapping. The failure-conditioned measurement
distribution is now characterized along the whole load axis (Proposition~\ref{prop:axis}, requiring $d/\sqrt M\to\infty$). The following questions are left for future work:
the error rate at finite size; in the moving subcritical region with $|r_M|/\sqrt{\log M}\to\infty$,
the relative probability and the fine conditional law of the macroscopic layer itself, a region beyond the scope of Proposition~\ref{prop:mda}, where the failure-conditioned law
is dominated by the repeated-edge layer; the complete classification of the sparse end $d=O(\sqrt M)$; $k$ growing with the size; and non-ideal checksums and other mapping families.
